% ═══════════════════════════════════════════════════════════════════
% CONCLUSION GÉNÉRALE
% ═══════════════════════════════════════════════════════════════════

\chapter*{Conclusion Générale}
\addcontentsline{toc}{chapter}{Conclusion Générale}

Ce rapport a rendu compte de l'ensemble du travail mené dans le cadre de ce stage d'été, de l'analyse des besoins à la validation d'une solution opérationnelle.

Le projet répondait à une problématique concrète de l'ARSII : la gestion dispersée et manuelle de ses contacts en Recherche et Innovation, qui empêchait une exploitation efficace de son réseau. Une plateforme web de gestion, d'organisation et d'enrichissement de bases de contacts a ainsi été analysée, conçue, réalisée, testée et déployée. Elle couvre l'ensemble des besoins identifiés : centralisation des contacts, automatisation de la saisie par import et numérisation OCR, exploitation fine du réseau par la recherche, la segmentation et le tableau de bord, et sécurisation de l'accès à deux rôles. Les tests fonctionnels et l'analyse statique ont validé le fonctionnement, et le déploiement cloud a rendu la plateforme accessible en ligne.

Sur le plan personnel, ce stage a permis d'acquérir une expérience concrète de développement d'une application web complète et du travail selon la méthodologie Agile/Scrum, de la spécification à la mise en production. La résolution d'une difficulté technique réelle (connexion Prisma/Supabase) a également développé la capacité d'investigation et de résolution de problèmes.

Plusieurs perspectives se dessinent : étendre la numérisation OCR à davantage de documents, affiner le modèle de rôles avec des permissions plus granulaires, et enrichir l'analyse du réseau par des indicateurs et visualisations supplémentaires dans le tableau de bord.
